% \chapter{Experiments}
% \label{cha:exper}
% 
% \section{Experiment Structure}
% 
% \section{Exploratory Analysis of Regret \& Runtime}
% \subsection{Experiment 1}
% \subsection{Experiment 2}
% \subsection{Experiment 3}
% \subsection{Experiment 4}
% 
% 
% \section{Expected Cumulative Regret and Runtime Profiles}
% % TODO: Clarify structure of such experiments
% % Gist: Math motiv -> setup -> collection of results (mentioned for both the experiment types in next chapter)
% \subsection{}

\chapter{Experiments}
\label{cha:exper}
 
\section{Experiment Structure}
\label{sec:exper:structure}
 
All experiments share a common infrastructure defined in the YAML configuration
files under \texttt{configs/experiments/} !TODO: Cite appendices!. 
Each configuration specifies four top-level blocks: \texttt{suite} 
(global settings), \texttt{problems} (benchmark functions), 
\texttt{algorithms} (optimisers), and \texttt{plotting} (output).
 
\subsection{Repetitions and budget}
Every algorithm--problem pair is evaluated over $R$ independent runs (typically
25--100, chosen to keep Monte-Carlo variance !CITE! in $\mathbb{E}[\text{CR}(T)]$
negligible), each run consuming up to a fixed evaluation budget $T_{\max}$
drawn from a per-suite budget list. A separate \texttt{budget\_for\_plots}
value, calibrated to approximately twice the expected time-to-first-optimum
($\mathbb{E}[\text{TTFO}]$) for the hardest algorithm on each problem, governs
the time horizon displayed in trajectory plots.
 
\subsection{Execution mode}
All suites run in \texttt{full} mode with parallelism enabled. 
The \texttt{full} supports recording the evaluation time-series. 
It was trivially toggled for these experiments.
Suites that compute inverse runtime profiles additionally set
\texttt{profile:~true}, which enables trajectory-level data collection required
for the tail-sum verification and runtime-profile plots.
 
% \subsection{Output}
% Raw results are written as \acrshort{json} under \texttt{results/raw/}, organized by suite,
% problem, and algorithm. Figures are written under \texttt{results/figures/}
% with a per-problem, per-size subdirectory structure~\ref{subsec:resint}. These results can be regenerated 
% with for each configuration of the experiment via the \acrshort{cli}.
% Plot types are described in Table~\ref{tab:plot-types}. (See Subsection~\ref{subsec:resint})

\subsection{Visualisation and Table Strategy}

The output of each suite consists of plots and summary tables generated from
the stored \acrshort{json} results. Plot types fall into four groups:
aggregate, budget-specific distribution, trajectory, and runtime profile.

\paragraph{Aggregate plots.}
Three plots pool results across all budgets in the suite budget list.
\emph{Mean simple-regret curves} (\texttt{regret\_curves}) plot
$\widehat{\mathbb{E}}[f^{*}-f(x_{\text{best}})]$ against evaluation budget
on a log--log axis with one shaded curve per algorithm; a steeper slope or
earlier crossover indicates a faster-decaying algorithm. When all algorithms
converge fully within the budget range this plot loses resolution, in which
case the \emph{convergence probability} plot (\texttt{convergence\_probability})
is more informative: it shows the fraction of runs reaching regret below
$10^{-9}$, separating algorithms by the budget at which they first converge
reliably. The \emph{comparison heatmap} (\texttt{comparison\_heatmap})
encodes $\log_{10}(\text{mean regret})$ as a colour over all
algorithm--budget cells, providing a compact global overview but at lower
precision than the regret curves.

\paragraph{Budget-specific distribution plots.}
These are generated once per problem at the budget specified by
\texttt{budget\_for\_plots}. \emph{Regret boxplots}
(\texttt{regret\_boxplots}) show the distribution of simple regret across
runs with individual points overlaid, exposing skewness and run-to-run
variability that means cannot: two algorithms can share the same mean while
one converges reliably on most runs and fails completely on a few.
\emph{Performance profiles} (\texttt{performance\_profile}) show the
empirical CDF of simple regret on a log-scaled axis; a curve that rises
steeply at low regret values indicates reliable near-optimal convergence,
whereas a gradual rise indicates high variability. Both plots support optional
pairwise statistical annotation (Mann--Whitney U and Cohen's $d$ against a
reference algorithm).

\paragraph{Trajectory plots.}
All trajectory plots are averaged across runs at \texttt{budget\_for\_plots},
with a standard-error band and optional TTFO markers at the mean
time-to-first-optimum. The suite produces several complementary variants.

The \emph{best-value history} (\texttt{history\_best}) plots the
monotonically non-decreasing incumbent fitness against evaluation time,
providing an absolute fitness view with $f^{*}$ as a horizontal reference.

\emph{Instantaneous regret} is produced in two variants.
With \texttt{track\_incumbent: false} (\texttt{regret\_instantaneous}),
it is the gap between $f^{*}$ and the solution evaluated at each step,
which fluctuates whenever an algorithm accepts a worsening move.
With \texttt{track\_incumbent: true} (\texttt{regret\_instantaneous\_best}),
it is the gap from $f^{*}$ to the incumbent, which is monotonically
non-increasing and equivalent to simple regret measured at each $t$.
For hill-climbing algorithms (RLS, $(1+1)$-EA) the two variants coincide
because no worsening moves are accepted; for SA and RLSExploration they
diverge during phases of active exploration. Both are plotted on a log-$y$
axis. A flat non-zero floor in the incumbent variant on this axis is the
primary diagnostic signature of saturation at a local optimum.

\emph{Cumulative regret} (\texttt{regret\_cumulative} and
\texttt{regret\_cumulative\_best}) is the integral of the corresponding
instantaneous regret series over time. It penalises slow convergence: an
algorithm that reaches $f^{*}$ late has accumulated all preceding regret
regardless of final solution quality. The slope of the cumulative regret curve
during the trapped phase grows linearly with time and is the key signal for
barrier-crossing suites (05, 06, 09, 10); this slope is suppressed on a
log axis, so cumulative regret is plotted on a linear $y$-axis throughout.

The \emph{TTFO distribution} (\texttt{ttfo\_distribution}) plots individual
run-level times-to-first-optimum as jittered points per algorithm. This plot
is conditioned on convergence: runs that did not reach the optimum contribute
no point, so when $P(\text{opt}) \approx 0$ across the suite the scatter is
sparse and adds little beyond the convergence probability plot.

\paragraph{Runtime profile plots.}
These require \texttt{suite.profile:~true} and are produced only in suites
03, 08, and 12. The inverse runtime profile
$P(\tau_v \leq T)$---the probability that fitness level $v$ was first reached
by evaluation $T$---is the common basis.

The \emph{profile surface} (\texttt{inverse\_runtime\_profile\_surface})
plots this probability as a heat map over $(T, v)$ for each algorithm
separately, with iso-probability contours overlaid. A front that stalls at
some $v_{\text{local}} < f^{*}$ regardless of additional budget identifies
the saturation level of the algorithm on that landscape.
The \emph{profile curves} (\texttt{inverse\_runtime\_profile\_curves}) plot
$P(\tau_v \leq T)$ versus $T$ at four selected fitness levels
($f^{*} \times \{0.25, 0.5, 0.75, 0.95\}$), overlaying all algorithms on
the same panel for direct comparison.

The \emph{profile verification} plot (\texttt{cr\_profile\_verification}),
enabled only in suite 08, overlays two independent estimates of
$\mathbb{E}[\text{CR}(T)]$: the empirical mean across runs, and the value
derived via the tail-sum identity from the inverse profile. For integer-valued
unit-increment fitness functions the identity is exact and any gap indicates a
computational error; for the continuous-valued NK landscape (suite 12) a minor
gap reflects finite-grid discretisation and not a bug.

\paragraph{Summary tables.}
Tables are exported per-suite in \LaTeX, CSV, or Markdown.
A \emph{per-problem table} is produced at each budget for each problem, with
one row per algorithm reporting mean, median, standard deviation, interquartile
range, a 95\% bootstrap confidence interval, $P(\text{opt})$, and run count,
sorted by ascending mean regret. The mean--median gap is informative when
distributions are skewed by a minority of non-converging runs.
An \emph{aggregate cross-problem table} summarises each algorithm's
across-problem mean regret, mean $P(\text{opt})$, and the count of problems
on which it achieved the lowest mean regret among all algorithms in the suite.


 
\subsection{Experiment numbering}
Suites are numbered 01--14. Suites 01--02 establish the baseline algorithm-family
comparison. Suites 03--07 and 09--14 form the exploratory analyses. Suite~08 is
a dedicated verification suite for the tail-sum identity.
Suites 03 and 12 additionally enable runtime profiling to study saturation
phenomena on NK landscapes alongside the verification suite.

\section{Runtime Profiles}


\subsection{Motivation}

The predominant performance measure in the runtime analysis of \acrshort{ea}s
is the \textit{expected optimization time}: the expected number of function evaluations
until an optimal solution is queried for the first
time~\cite{droste2002analysis,doerr2020theory}. While this measure has proved immensely
useful for ranking algorithms and establishing complexity hierarchies, it carries a
significant limitation for practice. In almost all realistic settings, a user cannot
detect when the algorithm has found an optimum. Even if this were possible, the
number of evaluations required to guarantee optimality may be prohibitively
large~\cite{pinto2018practice}. What the practitioner actually observes is the evolution
of solution quality over the \emph{entire} course of the run.

Two complementary frameworks have been proposed to address this limitation. The first,
from Jansen and Zarges~\cite{jansen2014fixed}, adopts a \textit{fixed-budget
perspective}: fix the evaluation budget $T$ in advance and characterize the expected
quality of the best solution found within that budget. The second, introduced by Pinto
and Doerr~\cite{pinto2018practice}, proposes \textit{runtime profiles}: rather than
reporting only the expected time to reach the global optimum, report the expected time
to reach each intermediate fitness level. These two views are orthogonal; one fixes
time and studies expected solution quality, the other fixes solution quality and studies
expected time. However, both provide richer information than the single expected
optimization time, as supported by Pinto and Doerr:
\begin{displayquote}
Our intention is rather to highlight to researchers working on theoretical aspects in
evolutionary computation that, beyond being possibly more relevant for practitioners,
explicitly reporting runtime profiles can give quite interesting insights into the
performance of \acrshort{ea}s.
{(\cite{pinto2018practice}, p.~19)}
\end{displayquote}

\subsection{Fixed-Budget Analysis}

In \textit{fixed-budget analysis}~\cite{jansen2014fixed}, the evaluation budget $T$ is
fixed in advance and the quantity of interest is the distribution of solution quality
achieved within that budget. Formally, let $f^{*}_T(\mathcal{A})$ denote the best
objective value achieved by algorithm $\mathcal{A}$ within $T$ evaluations on objective
$f$. Fixed-budget analysis characterizes $\mathbb{E}[f^{*}_T(\mathcal{A})]$ as a
function of $T$~\cite{jansen2014fixed}.

Fixed-budget analysis is particularly relevant in practice, where resources are finite
and an algorithm must return the best solution found when the budget is exhausted,
regardless of whether the global optimum has been identified~\cite{jansen2014fixed}. It
connects naturally to the simple regret framework: minimizing simple regret over a
fixed budget $T$ is equivalent to maximizing $\mathbb{E}[f^*_T(\mathcal{A})]$. Comparing
two algorithms under a fixed budget may yield different conclusions than comparing their
expected runtimes, since an algorithm with a larger expected optimization time may still
produce better intermediate solutions within a given budget~\cite{jansen2014fixed}.

\subsection{Variants}

\subsubsection{Expectation-Based Runtime Profiles}

Pinto and Doerr~\cite{pinto2018practice} propose that runtime results for evolutionary
algorithms should include not only the expected time to reach the global optimum but
also the expected time needed to reach each intermediate fitness level. Concretely, for
a fitness function $f : \{0,1\}^n \to \mathbb{R}$ with global optimum $f^* = \max_x
f(x)$, and for each target level $v$, the \textit{first hitting time} of level $v$ is
\begin{equation}
    \tau_v \;:=\; \min\bigl\{t \in \mathbb{N} : f(x_t) \geq v\bigr\},
\end{equation}
where $x_1, x_2, \ldots$ is the sequence of query points produced by the
algorithm~\cite{droste2002analysis}. The \textit{expectation-based runtime profile}
of an algorithm is the collection of expected first hitting times
$\bigl(\mathbb{E}[\tau_v]\bigr)_{v}$ across all relevant target
levels~\cite{pinto2018practice}.

When canonical fitness levels exist, such as for \textsc{OneMax} or
\textsc{LeadingOnes}, these provide natural targets; for other functions, intermediate
values identified during the analysis or obtained by linear interpolation of the minimal
and maximal fitness values may be used~\cite{pinto2018practice}. Crucially, Pinto and
Doerr emphasize that this framework does not require new algorithms or new proof
techniques, all bounds can be derived from existing runtime results and are often
already implicit in the proofs~\cite{pinto2018practice}. The proposal is therefore one
of \emph{presentation}: reporting richer, more informative performance summaries that
better reflect the anytime behavior of an algorithm.

The practical value of this richer view is illustrated by the behavior of (1+1)~\acrshort{ea}
variants on \textsc{LeadingOnes}: reporting only the expected optimization time suggests
RLS is superior, yet the runtime profile reveals that
(1+1)~EA$_{>0}$ (see !TODO: ref (1+1)-EA\_{>=0} variant!)
reaches all intermediate fitness levels $i \leq 7{,}980$ faster than RLS for problem
size $n = 10{,}000$, with RLS only pulling ahead in the final portion of the
run~\cite{pinto2018practice}. Such crossover behavior is invisible to the single
expected optimization time and is directly relevant to the design of adaptive parameter
and operator selection schemes~\cite{pinto2018practice}.

\subsubsection{Distributional Runtime Profiles}

The expectation-based runtime profile of~\cite{pinto2018practice} summarizes the
distribution of $\tau_v$ by its mean. A finer characterization, and the appropriate
quantity for connecting runtime profiles to regret-based measures, is the
\textit{distributional runtime profile}
\begin{equation}
    \mathcal{P}(v, t) \;:=\; \mathbb{P}(\tau_v > t),
\end{equation}
which records, for each target level $v$ and time step $t$, the probability that level
$v$ has not yet been reached by time $t$. The expectation-based profile is recovered by
the standard tail-sum formula for non-negative integer-valued random
variables~\cite{doerr2020theory}:
\begin{equation}
    \label{eqn:dist-runtime-rel}
    \mathbb{E}[\tau_v] \;=\; \sum_{t=0}^{\infty} \mathbb{P}(\tau_v > t).
\end{equation}
The distributional profile is the more primitive object from which expectation-based
statements follow, and it constitutes the link between runtime profiles and the
regret-based framework developed below.

\begin{remark}[Naming convention]
\label{rem:profile-naming}
The term \emph{runtime profile} appears in two distinct senses in the literature and in
this work. Pinto and Doerr~\cite{pinto2018practice} use it to denote the
\emph{expectation-based} runtime profile $\bigl(\mathbb{E}[\tau_v]\bigr)_v$, which
summarizes per-level performance by expected hitting times. The present work additionally
employs the strictly finer \emph{distributional} runtime profile $\mathcal{P}(v,t) =
\mathbb{P}(\tau_v > t)$, from which the expectation-based profile is recovered via the
tail-sum formula (Equation \ref{eqn:dist-runtime-rel}). To avoid ambiguity, the two notions are referred to explicitly whenever
both are in scope. Hereafter, and in all practical experimentation in this report,
\emph{runtime profile} refers to the distributional variant $\mathcal{P}(v,t)$ unless
otherwise stated.
\end{remark}

\subsection{Connection to Cumulative Regret}
\label{subsec:regret-profiles}

One of the discussions suggested by Pinto and Doerr in their \textit{Future Works}
section is highlighted:
\begin{displayquote}
    We feel that there is a need for combined measures that are capable of describing
    the anytime behavior of an EA[`Evolutionary Algorithms']. Regret-based measure as
    used in machine learning could be a key here, but we haven't been able so far to
    identify a fully satisfying measure.
    {(\cite{pinto2018practice}, p.~22)}
\end{displayquote}
As far as we know, establishing such a relationship remains an open question. This consolidated the following proposal for exploring the established concepts of \textit{\acrfull{cr}} and \textit{Runtime Profiles}.

The relationship between runtime profiles and \acrshort{cr} can be made precise
through a formal setup. A randomized black-box algorithm produces an adaptive sequence
of query points $x_1, x_2, \ldots \in \{0,1\}^n$. Define the
\textit{best-so-far process}
\[
    B(t) \;:=\; \max_{s \leq t} f(x_s),
\]
which records the best objective value seen up to time $t$. Since $B(t)$ is
non-decreasing in $t$, the event $\{B(t) < v\}$ is equivalent to $\{\tau_v > t\}$
(the level $v$ has not yet been hit). The \textit{\acrshort{cr}} with respect to the
global optimum $f^*$ over a budget of $T$ evaluations is~\cite{cesa2006prediction}
\[
    \mathrm{CR}(T) \;:=\; \sum_{t=1}^{T} \bigl(f^* - B(t)\bigr).
\]

The proof of the main theorem relies on two elementary lemmas.

\begin{lemma}[Tail-sum representation for the gap]
\label{lem:tailsum}
For any integer $b$ with $0 \le b \le f^*$,
\[
    f^* - b \;=\; \sum_{v=1}^{f^*} \mathbf{1}[b < v].
\]
\end{lemma}

\begin{proof}
The right-hand side counts the integers $v \in \{1,\ldots,f^*\}$ that strictly exceed
$b$, namely the set $\{b+1,\ldots,f^*\}$, which has cardinality $f^* - b$.
\end{proof}

\begin{lemma}[Indicator equivalence]
\label{lem:indicator}
For all $v \in \{1, \ldots, f^*\}$ and $t \in \mathbb{N}$,
\[
    \mathbf{1}[B(t) < v] \;=\; \mathbf{1}[\tau_v > t].
\]
\end{lemma}

\begin{proof}
Since $B(t)$ is non-decreasing in $t$, the event $\{B(t) \ge v\}$ is equivalent to
$\{\tau_v \le t\}$. Taking complements gives $\{B(t) < v\} = \{\tau_v > t\}$, so the
two indicators agree on every sample path.
\end{proof}

The following identity connects \textit{expected} cumulative regret to the distributional
runtime profile.

\begin{theorem}[Tail-Sum Regret Representation]
\label{thm:regret-profile}
For any randomized black-box algorithm operating on $f$ with finite budget
$T \in \mathbb{N}$,
\[
    \mathbb{E}[\mathrm{CR}(T)]
    \;=\;
    \sum_{v=1}^{f^*} \sum_{t=1}^{T} \mathbb{P}(\tau_v > t).
\]
\end{theorem}

\begin{proof}
For each $t$, Lemma~\ref{lem:tailsum} applied to $b = B(t)$ gives
\[
    f^* - B(t) \;=\; \sum_{v=1}^{f^*} \mathbf{1}[B(t) < v].
\]
Summing over $t = 1, \ldots, T$ yields
\[
    \mathrm{CR}(T)
    \;=\; \sum_{t=1}^{T}\bigl(f^* - B(t)\bigr)
    \;=\; \sum_{t=1}^{T} \sum_{v=1}^{f^*} \mathbf{1}[B(t) < v].
\]
By Lemma~\ref{lem:indicator}, $\mathbf{1}[B(t) < v] = \mathbf{1}[\tau_v > t]$ on every
sample path, so
\[
    \mathrm{CR}(T) \;=\; \sum_{t=1}^{T} \sum_{v=1}^{f^*} \mathbf{1}[\tau_v > t].
\]
Taking expectations and applying linearity of expectation together with
$\mathbb{E}[\mathbf{1}[\tau_v > t]] = \mathbb{P}(\tau_v > t)$ gives
\[
    \mathbb{E}[\mathrm{CR}(T)]
    \;=\; \sum_{t=1}^{T} \sum_{v=1}^{f^*} \mathbb{P}(\tau_v > t).
\]
Since both sums are finite, their order may be exchanged, yielding the stated result.
\end{proof}

Theorem~\ref{thm:regret-profile} expresses expected cumulative regret as a double sum
of the distributional runtime profile: $\mathbb{E}[\mathrm{CR}(T)]$ is the total
expected time the algorithm spends below each fitness level $v$, summed over all levels.
An algorithm with a uniformly smaller runtime profile (one that reaches each fitness
level earlier in expectation) incurs lower cumulative regret.

Taking the infinite-horizon limit, and applying the tail-sum identity for
$\mathbb{E}[\tau_v]$ along with the monotone convergence theorem, yields the following
asymptotic connection.

\begin{corollary}[Asymptotic Regret and Expected Hitting Times]
\label{cor:asymptotic-regret}
Suppose that $\tau_v$ is almost surely finite with finite expectation for each
$v \in \{1, \ldots, f^*\}$. Then
\[
    \lim_{T \to \infty} \mathbb{E}[\mathrm{CR}(T)]
    \;=\;
    \sum_{v=1}^{f^*} \mathbb{E}[\tau_v] \;-\; f^*.
\]
\end{corollary}
 
\begin{proof}
\textbf{Step 1: Tail-sum identity for $\mathbb{E}[\tau_v]$.}
Since $\tau_v \in \mathbb{N}_{\geq 1}$ by the convention that $t = 1$ is the first
oracle call and initialization is excluded from the query budget, we have
$\mathbb{P}(\tau_v > 0) = 1$. The standard tail-sum formula for a positive
integer-valued random variable then gives
\[
    \mathbb{E}[\tau_v]
    \;=\; \sum_{t=0}^{\infty} \mathbb{P}(\tau_v > t)
    \;=\; 1 + \sum_{t=1}^{\infty} \mathbb{P}(\tau_v > t),
\]
hence $\sum_{t=1}^{\infty} \mathbb{P}(\tau_v > t) = \mathbb{E}[\tau_v] - 1$.
 
\textbf{Step 2: Passing to the limit in Theorem~\ref{thm:regret-profile}.}
For each fixed $v$, the partial sums $\sum_{t=1}^{T}\mathbb{P}(\tau_v > t)$ are
non-decreasing in $T$ and non-negative, so the monotone convergence theorem gives
\[
    \lim_{T\to\infty} \sum_{t=1}^{T} \mathbb{P}(\tau_v > t)
    \;=\; \sum_{t=1}^{\infty} \mathbb{P}(\tau_v > t)
    \;=\; \mathbb{E}[\tau_v] - 1.
\]
 
\textbf{Step 3: Summing over $v$.}
Applying Step 2 to each $v = 1,\ldots, f^*$ and invoking
Theorem~\ref{thm:regret-profile},
\[
    \lim_{T \to \infty} \mathbb{E}[\mathrm{CR}(T)]
    \;=\; \sum_{v=1}^{f^*}\bigl(\mathbb{E}[\tau_v] - 1\bigr)
    \;=\; \sum_{v=1}^{f^*} \mathbb{E}[\tau_v] \;-\; f^*. \qedhere
\]
\end{proof}
 
Corollary~\ref{cor:asymptotic-regret} provides a unifying perspective: the long-run
expected cumulative regret decomposes as the sum of the expected first hitting times
over all fitness levels, minus a constant equal to $f^*$ (arising from the $t = 0$ term
in the tail-sum identity). Algorithms that reach all intermediate fitness levels quickly
are therefore precisely those that incur low long-run cumulative regret. This formalizes
the intuition, advanced in~\cite{pinto2018practice}, that an algorithm making steady
progress through the fitness landscape is preferable to one that stagnates at low
fitness values even if both ultimately find the optimum in the same expected time.
 
The identity also reveals that the regret-based measures studied in the online learning
literature~\cite{cesa2006prediction} and the runtime profile framework
of~\cite{pinto2018practice} are not independent approaches to measuring algorithm
performance but are two views of the same underlying quantity, connected through the
distributional runtime profile $\mathcal{P}(v, t)$.

\subsection{Practical Relevance}
\label{subsec:profiles-relevance}

The runtime profile framework addresses a fundamental limitation of the single expected
optimization time as a performance measure: it cannot distinguish between algorithms
that make steady progress and those that stagnate before eventually finding the optimum.
In practice, most applications operate under a finite evaluation budget, and
intermediate solution quality at every point in the run
matters~\cite{jansen2014fixed,pinto2018practice}.

Pinto and Doerr~\cite{pinto2018practice} note that runtime profiles are particularly
valuable for the design of non-static parameter and operator selection strategies in
EAs (such as adaptive mutation rates, crossover probabilities, and selection
pressure) since such strategies must make decisions based on current solution quality
rather than on the expected optimization time alone. The identification of
\emph{crossover fitness levels} (where one algorithm overtakes another in the runtime
profile) provides actionable guidance for when to switch between strategies.

Furthermore, as established in Theorem~\ref{thm:regret-profile}, any collection of
upper bounds on $\mathbb{E}[\tau_v]$ for all levels $v$ simultaneously yields an upper
bound on expected cumulative regret, and conversely, lower bounds on regret propagate to
lower bounds on a weighted aggregate of first hitting times. Cumulative regret thus
provides a single scalar summary of the entire runtime profile, complementing the
per-level view that the profile itself offers.

Pinto and Doerr~\cite{pinto2018practice} explicitly identify regret-based measures from
machine learning as a promising direction for capturing the anytime behavior of an
\acrshort{ea} within a unified complexity framework, while noting that identifying a
fully satisfying combined measure remains an open problem. The identity in
Theorem~\ref{thm:regret-profile} and Corollary~\ref{cor:asymptotic-regret} can be seen
as concrete progress toward that goal.


\section{Baseline Suites}
\subsection{Suite 01: Algorithm-Family Comparison}

Suite 01 compares four algorithm families across six pseudo-boolean problems
of increasing difficulty. The algorithms are: \acrshort{rls}, RLSExploration,
the $(1+1)$-\acrshort{ea}, and a single \acrshort{sa} variant with logarithmic cooling
$d=2.0$. Problems are OneMax, LeadingOnes, TwoMax, Jump-3, Trap-10, and BinVal,
all with $n=100$ except BinVal ($n=50$). Budgets range from 200 to 10\,000
across 50 independent runs.

BinVal is included for its exponential fitness structure but is excluded from
profile analysis: with $f^{*}=2^{50}-1$, constructing an integer fitness-level
grid is computationally infeasible. Profile plots are disabled suite-wide.
Pairwise statistical annotations are enabled on boxplots, performance profiles,
and \acrfull{ttfo} distributions, with \acrshort{rls} as the reference algorithm.

\subsection{Suite 02: Extended Pseudo-Boolean Benchmark}

Suite 02 extends the baseline to nine problems, adding Jump-5, Trap-5,
Plateau-10, HIFF-64, and a combinatorial instance (Petersen-3-Color).
The algorithm set is identical to suite 01 (the same four families), with
problem-specific hyperparameter overrides for RLSExploration ($\varepsilon=0.05$,
no decay on Jump-3, Jump-5, Trap-5) and SA-Log-d\_2.0 ($d=6.0$ on Trap-5 and
Jump-5). Budgets range from 500 to 20\,000 over 50 runs. Profile analysis is
disabled.



\section{Exploratory Analysis of Regret and Runtime}
 
The exploratory suites address four thematic research questions:
\begin{enumerate}
    \item Algorithm-family ordering across problem classes.
    \item The effect of landscape barrier structure and mutation rate on regret trajectory shape.
    \item The interaction between \acrshort{sa} cooling schedules and landscape ruggedness.
    \item Regret and convergence scaling with problem size and population structure.
\end{enumerate}
